\section{Conclusion}
\label{sec:conclusion}

We presented the first systematic study of ``taking offense'' as a behavioral and representational phenomenon in large language models.
Our multi-method design revealed three main findings.
First, LLMs process offense-relevant information and can articulate graded offense when permitted, but never express it spontaneously---an \offensegap that is largest for clearly offensive content and AI-directed existential provocations.
Second, a linear offense direction exists in hidden-state space that achieves perfect binary classification but fails to capture within-category severity, particularly for AI-directed prompts.
Third, the probe systematically diverges from behavioral offense, flagging directed criticism and existential statements as highly offensive even when the model reports minimal offense---a confound we trace to the probe detecting ``directed negative sentiment toward AI'' rather than offense per se.

Our convergent-validity framework provides practical guidance: offense probes can be cautiously trusted for binary screening (\AUROC $= 1.000$, zero false negatives) but should not be relied upon for fine-grained severity estimation or as training signals without behavioral validation.
More broadly, our results highlight a fundamental challenge for representation probing: when ground truth is contested, high classification accuracy is a necessary but insufficient condition for trust.

Future work should pursue causal interventions (activation patching at the identified offense layer), multi-model replication, human evaluation of the surprise cases, and control probes for related concepts (directed criticism, negative sentiment) to disentangle the offense signal from its confounds.